\documentclass{article}
\usepackage[utf8]{inputenc}
\usepackage{lscape}
\usepackage[normalem]{ulem}
\useunder{\uline}{\ul}{}


\title{Tarea 1}
\author{Alejandro Iabin Arteaga Hernández}
\date{Lenguajes de programación }




\usepackage{graphicx}

\begin{document}

\maketitle



% Please add the following required packages to your document preamble:
% \usepackage{booktabs}
% Please add the following required packages to your document preamble:
% \usepackage{booktabs}

\begin{landscape}
\begin{table}[]
\centering
\caption{Ejercicio 1}
\label{my-label}
\begin{tabular}{|l|l|l|l|l|l|l|}
\hline
\textit{\textbf{Nombre}} & \textit{\textbf{Año}} & \textit{\textbf{Paradigma}}                                   & \textit{\textbf{Diseñador}}                                 & \textit{\textbf{Usos principales}}                                                                                                            & \textit{\textbf{Principales características}}                                                                                                                                                                                                                                                                                                                                                                       & \textit{\textbf{Importancia en la historia}}                                                                                                                                                                                                                                                                                                                                                                                                                                                                                                                                                            \\ \hline
Php                      & 1995                  & \begin{tabular}[c]{@{}l@{}}Orientado\\ a objetos\end{tabular} & \begin{tabular}[c]{@{}l@{}}Rasmus\\ Ledorf\end{tabular}     & \begin{tabular}[c]{@{}l@{}}1.Scripting del lado \\ del servidor\\ 2.Desarrollo web\end{tabular}                                               & \begin{tabular}[c]{@{}l@{}}1.Cuenta con intérprete\\ 2.Completamente orientado\\ al desarrollo de aplicaciones\\ web dinámicas\\ 3.El código es invisible\\ al navegador, ya que todo\\ el código es ejecutado \\ por el servidor\\ 4.No requiere definición\\ de tipos de variables, ya\\ que el tipo se decide \\ en tiempo de ejecución\end{tabular}                                                             & \begin{tabular}[c]{@{}l@{}}El lenguaje mas usado del lado\\  del servidor, PHP es uno de \\ los pilares del internet moderno,\\ creado con la intención de hacer\\ páginas personales y dinámicas,\\ la importancia de PHP en la\\ actualidad es indiscutible.\\     Es un lenguaje altamente\\ robusto, usado por cualquier\\ tipo de usuario, desde páginas\\ personales, hasta sistemas de \\ una complejidad considerable. \\ Muchas páginas de gran\\ importancia fueron creadas con\\ PHP, por ejemplo: facebook.com,\\  wikipedia.org, etc.\end{tabular}                                         \\ \hline
C                        & 1972                  & Procedural                                                    & \begin{tabular}[c]{@{}l@{}}Dennis\\ Ritchie\end{tabular}    & \begin{tabular}[c]{@{}l@{}}En el desarrollo de:\\ 1.Sistemas operativos\\ 2.Sistemas embebidos\\ 3.Compiladores e \\ interpretes\end{tabular} & \begin{tabular}[c]{@{}l@{}}1.Lenguaje muy robusto,\\ con una gran cantidad de\\ funciones , que puede ser\\ usado para escribir \\ cualquier programa\\ complejo\\ 2.Rápido y eficiente\\ 3.Es altamente portable\\ 4.El compilador de C,\\ junta las capacidades\\  del lenguaje \\ ensamblador, con las \\ características de un\\ lenguaje de alto nivel\\ 5.Uso de punteros \\ altamente eficiente\end{tabular} & \begin{tabular}[c]{@{}l@{}}"El lenguaje de programación"\\ La importancia de C, no solo \\ abarca su valor histórico, sino \\ su valor en el mundo moderno.\\   La mayoría de los lenguajes \\ modernos, usan a C, no solo \\ para la programación de \\ compiladores e intérpretes, sino\\ también para la programación \\ de prácticamente todos los \\ sistemas operativos, los \\ programas de bajo nivel, que \\ se relacionan directamente con el\\ hardware. Usado como base de \\ mucha de la sintaxis de hoy en\\ día, C marcó una diferencia, para\\ el mundo de la informática.\end{tabular} \\ \hline
\end{tabular}
\end{table}
\end{landscape} 

\begin{landscape}
\begin{table}[]
\centering

\label{my-label}
\begin{tabular}{|l|l|l|l|l|l|l|}
\hline
\textit{\textbf{Nombre}} & \textit{\textbf{Año}} & \textit{\textbf{Paradigma}}                                   & \textit{\textbf{Diseñador}}                                 & \textit{\textbf{Usos principales}}                                                                                      & \textit{\textbf{Principales características}}                                                                                                                                                               & \textit{\textbf{Importancia en la historia}}                                                                                                                                                                                                                                                                                                                                                                                                                                                                                                                                      \\ \hline
C++                      & 1983                  & \begin{tabular}[c]{@{}l@{}}Orientado\\ a objetos\end{tabular} & \begin{tabular}[c]{@{}l@{}}Bjarne\\ Stroustrup\end{tabular} & \begin{tabular}[c]{@{}l@{}}1.Videojuegos\\ 2.Cómputo avanzado\\ 3.Gráficos\\ 4.Compiladores\end{tabular}                & \begin{tabular}[c]{@{}l@{}}1.Es una extensión \\ del lenguaje C\\ 2.Es muy rápido \\ y eficiente\\ 3.Estaticamente tipado\\ 4.Orientado a objetos\\ 5.Es compatible \\ con código escrito en C\end{tabular} & \begin{tabular}[c]{@{}l@{}}Creado originalmente como\\ una ampliación al lenguaje \\ C, incluyendo entre muchas\\ cosas, la orientación a objetos\\ C++, se convirtió por mucho\\ tiempo, en el lenguaje \\ predeterminado de aplicaciones\\ de escritorio, con el desarrollo \\ de nuevos lenguajes, ha perdido\\ popularidad, pero sigue \\ destacando por su gran eficiencia\\ frente a su competencia, siendo \\ el preferido para aplicaciones, \\ que necesiten un gran desempeño.\end{tabular}                                                                             \\ \hline
Python                   & 1991                  & \begin{tabular}[c]{@{}l@{}}Orientado\\ a objetos\end{tabular} & \begin{tabular}[c]{@{}l@{}}Guido\\ van Rossum\end{tabular}  & \begin{tabular}[c]{@{}l@{}}1.Aplicaciones web\\ 2.Inteligencia Artificial\\ 3.Aplicaciones de\\ escritorio\end{tabular} & \begin{tabular}[c]{@{}l@{}}1.Cuenta con un intérprete\\ 2.Fácil lectura\\ 3.Cuenta con un sistema de\\ tipos dinámico\\ 4.Soporta múltiples tipos\\ de paradigmas de \\ programación\end{tabular}           & \begin{tabular}[c]{@{}l@{}}La filosofía de python,\\ "un lenguaje accesible de leer"\\ cambió las reglas del juego de los \\ lenguajes actuales, creado con el\\ propósito de escribir código en\\ menos lineas que los lenguajes \\ tradicionales, y un sistema de \\ bloques sin uso de llaves, además\\ de funcionar como un lenguaje \\ multipropósito, que está presente\\ en muchos de los campos de la\\ informática moderna, desde\\ computo científico, hasta sistemas\\ de web scripting. Python está \\ cambiando la forma y evolución\\ de los lenguajes\end{tabular} \\ \hline
\end{tabular}
\end{table}
\end{landscape} 

\begin{landscape}

\begin{table}[]
\centering
\label{my-label}
\begin{tabular}{|l|l|l|l|l|l|l|}
\hline
\textit{\textbf{Nombre}} & \textit{\textbf{Año}} & \textit{\textbf{Paradigma}}                                   & \textit{\textbf{Diseñador}}                                                              & \textit{\textbf{Usos principales}}                                                                                            & \textit{\textbf{Principales características}}                                                                                                                                                                                             & \textit{\textbf{Importancia en la historia}}                                                                                                                                                                                                                                                                                                                                                                                         \\ \hline
JavaScript               & 1995                  & \begin{tabular}[c]{@{}l@{}}Orientado\\ a objetos\end{tabular} & \begin{tabular}[c]{@{}l@{}}Brendan \\ Eich\end{tabular}                                  & \begin{tabular}[c]{@{}l@{}}Frontend\\ scripting\end{tabular}                                                                  & \begin{tabular}[c]{@{}l@{}}1.Se ejecuta completamente\\ en el navegador\\ 2.Cuenta con una gran \\ cantidad de bibliotecas y \\ frameworks\\ 3.Dinamicamente tipado\end{tabular}                                                          & \begin{tabular}[c]{@{}l@{}}Uno de los grandes lenguajes\\ de la actualidad, junto con \\ HTML CSS forman los \\ pilares de la WWW, el \\ redescubrimiento de \\ javaScript permitió realizar\\ operaciones del lado del \\ cliente, cambiando \\ radicalmente lo que se podía\\ hacer en la web.\\     Actualmente es soportado\\ por todos los navegadores y\\ forma parte de la mayoría de\\ las páginas de internet.\end{tabular} \\ \hline
Haskell                  & 1990                  & Funcional                                                     & \begin{tabular}[c]{@{}l@{}}Universidad\\ de Yale\\ Universidad\\ de Glasgow\end{tabular} & Académico                                                                                                                     & \begin{tabular}[c]{@{}l@{}}1.Cuenta con evaluación\\ perezosa \\ 2.Semanticamente está \\ basado en el lenguaje\\ Miranda, pero con una \\ sintaxis enfocada a la \\ industria\\ 3.Cuenta con pattern \\ matching\end{tabular}            & \begin{tabular}[c]{@{}l@{}}Un lenguaje creado para la\\ academia, haskell ha \\ cambiado cosas, le ha \\ demostrado a la industria, que\\ la programación funcional\\ no es simplemente una \\ curiosidad científica más, y \\ ha introducido un nuevo \\ competidor a los lenguajes\\ comerciales. Empresas como\\ facebook, tienen\\ implementados grandes \\ sistemas en Haskell.\end{tabular}                                    \\ \hline
\end{tabular}
\end{table}
\end{landscape}

\begin{landscape}
\begin{table}[]
\centering

\label{my-label}
\begin{tabular}{|l|l|l|l|l|l|l|}
\hline
\textit{\textbf{Nombre}} & \textit{\textbf{Año}} & \textit{\textbf{Paradigma}}                                   & \textit{\textbf{Diseñador}}                                 & \textit{\textbf{Usos principales}}                                                                                            & \textit{\textbf{Principales características}}                                                                                                                                                                                             & \textit{\textbf{Importancia en la historia}}                                                                                                                                                                                                                                                                         \\ \hline
C\#                      & 2000                  & \begin{tabular}[c]{@{}l@{}}Orientado\\ a objetos\end{tabular} & Microsoft                                                   & \begin{tabular}[c]{@{}l@{}}1.Aplicaciones de\\ escritorio\\ 2.Desarrollo web\end{tabular}                                     & \begin{tabular}[c]{@{}l@{}}1.Se enfoca mucho en la\\ eficiencia del programador\\ 2.Arquitectura orientada a\\ objetos totalmente\\ 3.Cuenta con todo el soporte\\ de Microsoft, y es \\ compatible con la plataforma\\ .NET\end{tabular} & \begin{tabular}[c]{@{}l@{}}Pieza fundamental en la \\ estrategia de monopolizar los \\ servicios de cómputo por \\ parte de Microsoft, C\# es la \\ apuesta, lenguaje fundamental\\ en las herramientas de microsoft\\ C\# es uno de los grandes \\ lenguajes de hoy en día\end{tabular}                             \\ \hline
Perl                     & 1987                  & Funcional                                                     & Larry Wall                                                  & \begin{tabular}[c]{@{}l@{}}1.Administración\\ de sistemas\\ 2.network \\ programming\\ 3.Análisis \\ estadístico\end{tabular} & \begin{tabular}[c]{@{}l@{}}1.Lenguaje interpretado\\ 2.Fue creado como un \\ lenguaje de scripting para \\ unix\end{tabular}                                                                                                              & \begin{tabular}[c]{@{}l@{}}Sobresalió por su facilidad de\\ manejo de expresiones regulares\\ pero su punto fuerte, siempre fue\\ la administración de sistemas,\\ empresas como amazon usan \\ perl.\end{tabular}                                                                                                   \\ \hline

\end{tabular}
\end{table}
\end{landscape} 

\begin{landscape}

\begin{table}[]
\centering

\label{my-label}
\begin{tabular}{|l|l|l|l|l|l|l|}
\hline
\textit{\textbf{Nombre}} & \textit{\textbf{Año}} & \textit{\textbf{Paradigma}} & \textit{\textbf{Diseñador}}                                 & \textit{\textbf{Usos principales}}                                                                   & \textit{\textbf{Principales características}}                                                                                                                      & \textit{\textbf{Importancia en la historia}}                                                                                                                                                                                                                                                                         \\ \hline
Scratch                  & 2005                  & Event-driven                & MIT                                                         & Enseñanza                                                                                            & \begin{tabular}[c]{@{}l@{}}1.Es un lenguaje visual\\ 2.Es Turing-completo\\ 3.Está enfocado a la\\ enseñanza\end{tabular}                                          & \begin{tabular}[c]{@{}l@{}}Scratch, es un lenguaje de \\ programación, completamente\\ orientado a la enseñanza, \\ bajo una filosofía, de que todos\\ debería saber programar, su \\ enfoque son los niños, Scratch es\\ la representación de esta nueva \\ filosofía, que cada vez toma \\ más fuerza\end{tabular} \\ \hline
Prolog                   & 1972                  & Lógico                      & \begin{tabular}[c]{@{}l@{}}Alain\\  Colmerauer\end{tabular} & \begin{tabular}[c]{@{}l@{}}1.Inteligencia\\ artificial\\ 2.Lingüistica \\ computacional\end{tabular} & \begin{tabular}[c]{@{}l@{}}1.Lenguaje declarativo\\ 2.Se basa en un sistema de \\ hechos y reglas\\ 3.implementa procesamiento \\ de lenguaje natural\end{tabular} & \begin{tabular}[c]{@{}l@{}}Fue uno de los primeros \\ lenguajes lógicos, dando así\\ nacimiento a muchísimas ramas\\ de la CS , fue sobresaliente en el\\ desarrollo de inteligencia \\ artificial por su sistema de \\ deducción natural.\end{tabular}                                                              \\ \hline
\end{tabular}
\end{table}
\end{landscape} 

\textbf{1.2}\\
\textbf{\textit{1. ¿Cuál es el nombre del lenguaje?}}
\\\\
Java\\\\
\textbf{\textit{2. ¿Cuál es el nombre o nombres de los diseñadores del lenguaje?}}
 \\\\
James Gosling\\\\
\textbf{\textit{3. ¿Cuál es el año de creación del lenguaje?}}
\\\\
1995\\\\
\textbf{\textit{4. ¿A qué paradigma pertenece el lenguaje de programación?}}
\\\\
Orientado a Objetos\\\\
\textbf{\textit{5. ¿El lenguaje tiene características de otro paradigma? Si la respuesta es afirmativa, mencionar dichas características.}}\\

Sí, Java por cuestión de diseño, no es completamente orientado a objetos, ya que no todos los valores son objetos, hay tipos básicos, que están fuera de toda definición de objeto, heredado de la programación imperativa.\\\\
\textbf{\textit{6. ¿Cuál es el uso principal del lenguaje?}}
\\\\
La creación de aplicaciones android\\\\
\textbf{\textit{7. ¿Cuáles son las características principales de este lenguaje?}}\\
Es un lenguaje compilado, traduce el código a byteCode, que es interpretado y ejecutado en una máquina virtual. Esta máquina virtual, la JVM, es la principal virtud de Java, que le permite ser ejecutado prácticamente en cualquier arquitectura. Otra característica muy importante es la orientación a objetos, la cual está sintácticamente obligada.\\\\
\textbf{\textit{8. ¿Con qué propósito fue creado el lenguaje?}}
\\\\
Fue creado con la intención de crear un lenguaje universal, que fuese capaz de correr sobre cualquier arquitectura, sin modificar el código.
\\\\
\begin{thebibliography}{9}

\bibitem{C} 
Wikipedia: The C Programming Language   18 ago 2017
\\\texttt{https://en.wikipedia.org/wiki/The\_C\_Programming\_Language}


\bibitem{C++} 
Wikipedia: C++   18 ago 2017
\\\texttt{https://en.wikipedia.org/wiki/C\%2B\%2B}

\bibitem{php} 
Wikipedia: PHP   18 ago 2017
\\\texttt{https://en.wikipedia.org/wiki/PHP}

\bibitem{C#} 
Wikipedia: C Sharp (programming language)   18 ago 2017
\\\texttt{https://en.wikipedia.org/wiki/C\_Sharp\_(programming\_language)}

\bibitem{prolog} 
Wikipedia: Prolog   18 ago 2017
\\\texttt{https://en.wikipedia.org/wiki/Prolog}

\bibitem{scratch} 
Wikipedia: Scratch Programming Language   18 ago 2017
\\\texttt{https://en.wikipedia.org/wiki/Scratch\_(programming\_language)}

\bibitem{haskell} 
Wikipedia: Haskell Programming Language   18 ago 2017
\\\texttt{https://en.wikipedia.org/wiki/Haskell\_(programming\_language)}

\bibitem{JavaScript} 
Wikipedia: JavaScript   18 ago 2017
\\\texttt{https://en.wikipedia.org/wiki/JavaScript}

\bibitem{perl} 
Wikipedia: Perl  18 ago 2017
\\\texttt{https://en.wikipedia.org/wiki/Perl}

\bibitem{python} 
Wikipedia: Python Programming Language   18 ago 2017
\\\texttt{https://en.wikipedia.org/wiki/Python\_(programming\_language)}

\bibitem{java} 
Wikipedia: Java Programming Language   18 ago 2017
\\\texttt{https://en.wikipedia.org/wiki/Java\_(programming\_language))}

\bibitem{java} 
normasapa.com: Como citar-referenciar wikipedia normas apa \\ 18 ago 2017
\\\texttt{http://normasapa.com/como-citar-referenciar-wikipedia-normas-apa/}




\end{thebibliography}
\end{document}
